Write
\[
 \phi(u)=Ae^{-\beta u}\mathbf 1_{\{u>0\}},\qquad
 d\widetilde N(t)=dN(t)-\lambda\,dt,\qquad
 dM(t)=dN(t)-\ell(t)\,dt.
\tag{1}
\]
Fix a law in \(\mathcal M_{\Delta,p}^{\mathrm{stab}}\).  By \(\text{def:stable-common-exponential-class}\), \(\Delta>0\), \(\mu_i>0\), \(\beta>0\), \(A\) is entrywise nonnegative, and \(\rho(A/\beta)<1\).  Taking expectations in the nonnegative intensity equation is justified componentwise by monotone convergence.  Stationarity then gives
\[
 \lambda=\mu+\int_0^\infty Ae^{-\beta u}\lambda\,du
 =\mu+(A/\beta)\lambda.
\tag{2}
\]
In particular \(\lambda_i\geq\mu_i>0\) for every \(i\), so
\(D=\operatorname{diag}(\lambda)\succ0\).

If \(z\in\operatorname{spec}(A)\), then \(|z|\leq\rho(A)<\beta\).  Therefore every eigenvalue \(z-\beta\) of \(B=A-\beta I_p\) satisfies
\[
 \operatorname{Re}(z-\beta)\leq |z|-\beta<0.
\tag{3}
\]
Thus \(B\) is a real Hurwitz matrix.  In particular, there are constants \(c>0\) and \(C<\infty\) such that \(\|e^{Bu}\|\leq Ce^{-cu}\) for \(u\geq0\), after enlarging \(C\) to absorb the polynomial factors of the finitely many Jordan blocks.  Every improper matrix integral used below therefore converges absolutely.

Subtracting (2) from the intensity equation and using the definition of the compensated martingale measure in (1) gives the centered innovation equation
\[
 d\widetilde N=dM+\phi*d\widetilde N.
\tag{4}
\]
For every integer \(d\geq1\), induction and the beta integral give
\[
 \phi^{*d}(u)
 =A^d\frac{u^{d-1}}{(d-1)!}e^{-\beta u}\mathbf 1_{\{u>0\}}.
\tag{5}
\]
Moreover, \(\int_0^\infty\phi^{*d}(u)\,du=(A/\beta)^d\) entrywise.  Since \(A/\beta\geq0\) and its spectral radius is less than one, the matrix geometric series converges.  Hence the entrywise \(L^1\) resolvent exists and (5) sums to
\[
 \psi(u):=\sum_{d=1}^{\infty}\phi^{*d}(u)
 =e^{-\beta u}e^{Au}A=e^{Bu}A,\quad u>0.
\tag{6}
\]
Iteration of (4) therefore yields, on every bounded interval in second mean,
\[
 d\widetilde N=dM+\psi*dM.
\tag{7}
\]
Indeed, the predictable quadratic covariation of \(M\) is \(D\,dt\), so the compensated-integral isometry reduces convergence of the remainder to the entrywise \(L^1\) convergence above together with the finite intensity in (2).

We next compute the positive-lag covariance with the orientation in the statement: the future coordinate is the row and the time-zero coordinate is the column.  For \(t>0\), causality in (7), orthogonality of martingale increments on disjoint intervals, and the predictable covariance measure \(D\,dt\) give
\[
 \begin{aligned}
 C_{+}(t)
 &=\psi(t)D+
   \int_{-\infty}^{0}\psi(t-s)D\psi(-s)^{\mathsf T}\,ds\\
 &=\psi(t)D+int_0^\infty\psi(t+u)D\psi(u)^{\mathsf T}\,du.
 \end{aligned}
\tag{8}
\]
Simplicity supplies the separate diagonal atom \(D\delta_0\); it does not enter (8), which is the nonsingular density at \(t>0\).  Substituting (6), and using that \(A\), \(B\), and \(e^{Bu}\) commute, yields
\[
 C_{+}(t)
 =e^{Bt}\left\{AD+AQA^{\mathsf T}\right\},\quad
 Q:=\int_0^\infty e^{Bu}De^{B^{\mathsf T}u}\,du.
\tag{9}
\]
For every nonzero \(x\in\mathbb R^p\),
\[
 x^{\mathsf T}Qx
 =\int_0^\infty
 \left\|D^{1/2}e^{B^{\mathsf T}u}x\right\|_2^2,du>0,
\tag{10}
\]
because the continuous integrand at \(u=0\) equals \(\|D^{1/2}x\|_2^2>0\).  Thus \(Q\succ0\).  Differentiating the integrand defining \(Q\), integrating over \([0,\infty)\), and using the Hurwitz limit gives
\[
 BQ+QB^{\mathsf T}=-D,\quad
 D=2\beta Q-AQ-QA^{\mathsf T}.
\tag{11}
\]
Consequently the bracket in (9) is
\[
 \begin{aligned}
 K&:=AD+AQA^{\mathsf T}\\
 &=A(2\beta Q-AQ-QA^{\mathsf T})+AQA^{\mathsf T}\\
 &=A(2\beta I_p-A)Q,
 \end{aligned}
\tag{12}
\]
and (9) proves \(C_{+}(t)=e^{Bt}K\).  Since \(AQA^{\mathsf T}\) is symmetric, the first expression in (12) also gives
\[
 K-K^{\mathsf T}=AD-DA^{\mathsf T}.
\tag{13}
\]

Fix \(k\geq1\).  The bins \([0,\Delta)\) and \([k\Delta,(k+1)\Delta)\) are disjoint; for \(k=1\) their common endpoint has zero covariance measure.  Parametrize the later time by \(k\Delta+u\) and the earlier time by \(v\), with \(u,v\in[0,\Delta]\).  By \(\text{def:population-moments}\), stationarity, (9), and Fubini's theorem,
\[
 \begin{aligned}
 \Gamma_k
 &=\int_0^\Delta\!\int_0^\Delta
 e^{B(k\Delta+u-v)}K\,du\,dv\\
 &=e^{B(k-1)\Delta}
 \left(\int_0^\Delta e^{Bu}\,du\right)
 \left(\int_0^\Delta e^{B(\Delta-v)}\,dv\right)K\\
 &=T^{k-1}J^2K.
 \end{aligned}
\tag{14}
\]
All uses of Fubini are valid because the integrands are continuous on the compact bin square.

It remains to prove the image assertions.  Since \(B\) is Hurwitz, it is invertible.  Spectral mapping and (3) show that every eigenvalue of \(T=e^{B\Delta}\) has modulus strictly less than one; hence \(T-I_p\) is invertible.  Therefore
\[
 J=\int_0^\Delta e^{Bu}\,du=B^{-1}(T-I_p)
\tag{15}
\]
is invertible.  Also \(2\beta I_p-A\) is invertible: if \(z\in\operatorname{spec}(A)\), then \(|z|<\beta<2\beta\).  Since \(Q\) is invertible by (10), (12) implies
\[
 \operatorname{im}(K)=\operatorname{im}(A).
\tag{16}
\]
The matrices \(J\) and \(A\) commute because \(J\) is an integral of functions of \(B=A-\beta I_p\); their inverses commute as well.  Hence \(J^2\) maps \(\operatorname{im}(A)\) bijectively onto itself.  Equations (14) at \(k=1\) and (16) now yield
\[
 \operatorname{im}(\Gamma_1)=J^2\operatorname{im}(K)
 =J^2\operatorname{im}(A)=\operatorname{im}(A).
\tag{17}
\]
Together, (16)--(17) prove \(\operatorname{im}(\Gamma_1)=\operatorname{im}(K)=\operatorname{im}(A)\), and taking dimensions proves \(\operatorname{rank}(\Gamma_1)=\operatorname{rank}(A)\).